% GENERATED FILE. Edit canonical lessons, then run npm run generate:books.
% canonical-source-hash: fnv1a64:186b13ea84608dc9
% canonical-lessons: TE-C337-arestu-ceyu, TE-C337-pattubadu, TE-C337-khaidi, TE-C337-sankellu, TE-C337-ninditudu

\chapter{Under Arrest}
\label{ch:te-under-arrest}
\hlchaptermodality{337}

\begin{chapteropening}
\textbf{By the end of this chapter:} \emph{I can say to arrest, to get caught, a prisoner, shackles and the accused.}

Complete the last lesson of chapter 337: I can say to arrest, to get caught, a prisoner, shackles and the accused.
\end{chapteropening}

\section[\texorpdfstring{\te{అరెస్టు} \te{చేయు} (aresṭu cēyu)}{అరెస్టు చేయు (aresṭu cēyu)}]{\te{అరెస్టు} \te{చేయు} (aresṭu cēyu) — to arrest}

\label{lesson:TE-C337-arestu-ceyu}



\begin{quote}
Before the new one: say the Telugu for an uproar, then the Telugu for to disclose.
\end{quote}

\subsection*{You'll want to know: \te{అరెస్టు} \te{చేయు}}
\textbf{\te{అరెస్టు} \te{చేయు}} — \emph{aresṭu cēyu} — \textquotedblleft{}to arrest\textquotedblright{}.

\begin{quote}
\te{పోలీసులు} \te{దొంగను} \te{అరెస్టు} \te{చేశారు}.
\end{quote}

\begin{quote}
\emph{pōlīsulu doṅganu aresṭu cēśāru.}
\end{quote}

\begin{quote}
The police arrested the thief.
\end{quote}

Past (I …): \te{అరెస్టు} \te{చేశాను}, \emph{aresṭu cēśānu}.

An English word with \te{చేయు}, as in \te{టైప్} \te{చేయు}; who is arrested takes -\te{ను}.

\begin{cousinweb}
\textbf{\te{అరెస్టు} \te{చేయు}} (\emph{aresṭu cēyu}) is a loan from English \emph{arrest} with the native verb \emph{cēyu}, \textquotedblleft{}to do\textquotedblright{}.
\end{cousinweb}

\begin{grammarlens}[title={What you've built}]
One more word for telling what happened.
\end{grammarlens}

\subsection*{Your turn}
\begin{itemize}
\raggedright
  \item \emph{Say it:} \emph{aresṭu cēyu}
  \item \emph{Say it:} \emph{aresṭu cēyu}, once more
  \item \emph{Say it:} \emph{pōlīsulu doṅganu aresṭu cēśāru.} — The police arrested the thief.
  \item \emph{Recall:} say the Telugu for an uproar, then the Telugu for to disclose, then say \emph{aresṭu cēyu} again
\end{itemize}

\begin{tcolorbox}[breakable,colback=teal!4,colframe=teal!35!black,title={Before you move on}]
What does \te{అరెస్టు} \te{చేయు} mean? (\textquotedblleft{}To arrest\textquotedblright{}.) Say it once more.
\end{tcolorbox}
\section[\texorpdfstring{\te{పట్టుబడు} (paṭṭubaḍu)}{పట్టుబడు (paṭṭubaḍu)}]{\te{పట్టుబడు} (paṭṭubaḍu) — to get caught}

\label{lesson:TE-C337-pattubadu}



\begin{quote}
Before the new one: say the Telugu for to disclose, then the Telugu for to arrest.
\end{quote}

\subsection*{You'll want to know: \te{పట్టుబడు}}
\textbf{\te{పట్టుబడు}} — \emph{paṭṭubaḍu} — \textquotedblleft{}to get caught\textquotedblright{}.

\begin{quote}
\te{దొంగ} \te{బస్టాండులో} \te{పట్టుబడ్డాడు}.
\end{quote}

\begin{quote}
\emph{doṅga basṭāṇḍulō paṭṭubaḍḍāḍu.}
\end{quote}

\begin{quote}
The thief got caught at the bus stand.
\end{quote}

Past (I …): \te{పట్టుబడ్డాను}, \emph{paṭṭubaḍḍānu}.

\te{పట్టు}, to catch, with \te{పడు}, which turns an action into something that happens to you.

\begin{cousinweb}
\textbf{\te{పట్టుబడు}} (\emph{paṭṭubaḍu}) is a native compound: \emph{paṭṭu}, \textquotedblleft{}to hold, catch\textquotedblright{} (DEDR 4034), with \emph{paḍu}, the helping verb that makes it something suffered.
\end{cousinweb}

\begin{grammarlens}[title={What you've built}]
One more word for telling what happened.
\end{grammarlens}

\subsection*{Your turn}
\begin{itemize}
\raggedright
  \item \emph{Say it:} \emph{paṭṭubaḍu}
  \item \emph{Say it:} \emph{paṭṭubaḍu}, once more
  \item \emph{Say it:} \emph{doṅga basṭāṇḍulō paṭṭubaḍḍāḍu.} — The thief got caught at the bus stand.
  \item \emph{Recall:} say the Telugu for to disclose, then the Telugu for to arrest, then say \emph{paṭṭubaḍu} again
\end{itemize}

\begin{tcolorbox}[breakable,colback=teal!4,colframe=teal!35!black,title={Before you move on}]
What does \te{పట్టుబడు} mean? (\textquotedblleft{}To get caught\textquotedblright{}.) Say it once more.
\end{tcolorbox}
\section[\texorpdfstring{\te{ఖైదీ} (khaidī)}{ఖైదీ (khaidī)}]{\te{ఖైదీ} (khaidī) — a prisoner}

\label{lesson:TE-C337-khaidi}



\begin{quote}
Before the new one: say the Telugu for to arrest, then the Telugu for to get caught.
\end{quote}

\subsection*{You'll want to know: \te{ఖైదీ}}
\textbf{\te{ఖైదీ}} — \emph{khaidī} — \textquotedblleft{}a prisoner\textquotedblright{}.

\begin{quote}
\te{ఖైదీ} \te{జైలు} \te{నుంచి} \te{పారిపోయాడు}.
\end{quote}

\begin{quote}
\emph{khaidī jailu nuñci pāripōyāḍu.}
\end{quote}

\begin{quote}
The prisoner escaped from jail.
\end{quote}

Plural: \te{ఖైదీలు}, \emph{khaidīlu}.

From \te{ఖైదు}, prison or custody: \te{ఖైదీ} is the one held.

\begin{cousinweb}
\textbf{\te{ఖైదీ}} (\emph{khaidī}) is a loan from Persian \emph{qaidī}, \textquotedblleft{}a prisoner\textquotedblright{}.
\end{cousinweb}

\begin{grammarlens}[title={What you've built}]
One more word for telling what happened.
\end{grammarlens}

\subsection*{Your turn}
\begin{itemize}
\raggedright
  \item \emph{Say it:} \emph{khaidī}
  \item \emph{Say it:} \emph{khaidī}, once more
  \item \emph{Say it:} \emph{khaidī jailu nuñci pāripōyāḍu.} — The prisoner escaped from jail.
  \item \emph{Recall:} say the Telugu for to arrest, then the Telugu for to get caught, then say \emph{khaidī} again
\end{itemize}

\begin{tcolorbox}[breakable,colback=teal!4,colframe=teal!35!black,title={Before you move on}]
What does \te{ఖైదీ} mean? (\textquotedblleft{}A prisoner\textquotedblright{}.) Say it once more.
\end{tcolorbox}
\section[\texorpdfstring{\te{సంకెళ్ళు} (saṅkeḷḷu)}{సంకెళ్ళు (saṅkeḷḷu)}]{\te{సంకెళ్ళు} (saṅkeḷḷu) — shackles; chains}

\label{lesson:TE-C337-sankellu}



\begin{quote}
Before the new one: say the Telugu for to get caught, then the Telugu for a prisoner.
\end{quote}

\subsection*{You'll want to know: \te{సంకెళ్ళు}}
\textbf{\te{సంకెళ్ళు}} — \emph{saṅkeḷḷu} — \textquotedblleft{}shackles; chains\textquotedblright{}.

\begin{quote}
\te{పోలీసులు} \te{అతని} \te{చేతులకు} \te{సంకెళ్ళు} \te{వేశారు}.
\end{quote}

\begin{quote}
\emph{pōlīsulu atani cētulaku saṅkeḷḷu vēśāru.}
\end{quote}

\begin{quote}
The police put chains on his hands.
\end{quote}

A plural word: \te{సంకెళ్ళు} \te{వేయు}, to put in chains.

\begin{cousinweb}
\textbf{\te{సంకెళ్ళు}} (\emph{saṅkeḷḷu}) is usually traced to Sanskrit \emph{śṛṅkhalā}, \textquotedblleft{}a chain, a fetter\textquotedblright{} (Monier-Williams).
\end{cousinweb}

\begin{grammarlens}[title={What you've built}]
One more word for telling what happened.
\end{grammarlens}

\subsection*{Your turn}
\begin{itemize}
\raggedright
  \item \emph{Say it:} \emph{saṅkeḷḷu}
  \item \emph{Say it:} \emph{saṅkeḷḷu}, once more
  \item \emph{Say it:} \emph{pōlīsulu atani cētulaku saṅkeḷḷu vēśāru.} — The police put chains on his hands.
  \item \emph{Recall:} say the Telugu for to get caught, then the Telugu for a prisoner, then say \emph{saṅkeḷḷu} again
\end{itemize}

\begin{tcolorbox}[breakable,colback=teal!4,colframe=teal!35!black,title={Before you move on}]
What does \te{సంకెళ్ళు} mean? (\textquotedblleft{}Shackles; chains\textquotedblright{}.) Say it once more.
\end{tcolorbox}
\section[\texorpdfstring{\te{నిందితుడు} (nindituḍu)}{నిందితుడు (nindituḍu)}]{\te{నిందితుడు} (nindituḍu) — the accused; a suspect}

\label{lesson:TE-C337-ninditudu}



\begin{quote}
Before the new one: say the Telugu for a prisoner, then the Telugu for shackles.
\end{quote}

\subsection*{You'll want to know: \te{నిందితుడు}}
\textbf{\te{నిందితుడు}} — \emph{nindituḍu} — \textquotedblleft{}the accused; a suspect\textquotedblright{}.

\begin{quote}
\te{నిందితుడు} \te{తాను} \te{దొంగ} \te{కాదని} \te{చెప్పాడు}.
\end{quote}

\begin{quote}
\emph{nindituḍu tānu doṅga kādani ceppāḍu.}
\end{quote}

\begin{quote}
The accused said that he was not the thief.
\end{quote}

Built on the Sanskrit word for blame, \te{నింద}; a woman is \te{నిందితురాలు}.

\begin{cousinweb}
\textbf{\te{నిందితుడు}} (\emph{nindituḍu}) is borrowed from Sanskrit \emph{nindita}, \textquotedblleft{}blamed, censured\textquotedblright{} (CDIAL 7213), with the Telugu masculine ending \emph{-uḍu}.
\end{cousinweb}

\begin{grammarlens}[title={What you've built}]
One more word for telling what happened.
\end{grammarlens}

\subsection*{Your turn}
\begin{itemize}
\raggedright
  \item \emph{Say it:} \emph{nindituḍu}
  \item \emph{Say it:} \emph{nindituḍu}, once more
  \item \emph{Say it:} \emph{nindituḍu tānu doṅga kādani ceppāḍu.} — The accused said that he was not the thief.
  \item \emph{Recall:} say the Telugu for a prisoner, then the Telugu for shackles, then say \emph{nindituḍu} again
\end{itemize}

\begin{tcolorbox}[breakable,colback=teal!4,colframe=teal!35!black,title={Before you move on}]
What does \te{నిందితుడు} mean? (\textquotedblleft{}The accused; a suspect\textquotedblright{}.) Say it once more.
\end{tcolorbox}
